\documentclass[aspectratio=169,11pt]{beamer}
\usepackage{fontspec}
\usepackage{listings}
\usepackage{xcolor}
\usepackage{booktabs}
\usepackage{array}

% ================= Palette (identique au livre Tome 1) =================
\definecolor{PyBleu}{HTML}{1A4A7A}
\definecolor{PyBleuMoyen}{HTML}{2D7AB3}
\definecolor{PyBleuClair}{HTML}{EAF1F8}
\definecolor{PyJaune}{HTML}{FFD43B}
\definecolor{PyNuit}{HTML}{142840}
\definecolor{PyTexte}{HTML}{1F2A37}
\definecolor{PyGris}{HTML}{5A6B7E}
\definecolor{PyRouge}{HTML}{B3442D}
\definecolor{PyVert}{HTML}{2E7D5B}
\definecolor{CodeTexte}{HTML}{E8EEF6}

% ================= Code Python =================
\lstdefinestyle{py}{
  language=Python,
  keywordstyle=\color{PyJaune}\bfseries,
  stringstyle=\color{PyJaune},
  commentstyle=\color{PyGris}\itshape,
  showstringspaces=false,
  breaklines=true,
  backgroundcolor=\color{PyNuit},
  basicstyle=\ttfamily\footnotesize\color{CodeTexte},
  columns=fullflexible,
  keepspaces=true,
  xleftmargin=6pt,
  framexleftmargin=6pt,
}
\lstdefinestyle{pysmall}{
  language=Python,
  keywordstyle=\color{PyJaune}\bfseries,
  stringstyle=\color{PyJaune},
  commentstyle=\color{PyGris}\itshape,
  showstringspaces=false,
  breaklines=true,
  backgroundcolor=\color{PyNuit},
  basicstyle=\ttfamily\scriptsize\color{CodeTexte},
  columns=fullflexible,
  keepspaces=true,
  xleftmargin=5pt,
  framexleftmargin=5pt,
}

% ================= Titres de diapositives =================
\setbeamercolor{structure}{fg=PyBleu}
\setbeamerfont{frametitle}{series=\bfseries, size=\Large}
\setbeamercolor{frametitle}{bg=PyBleu, fg=white}
\setbeamertemplate{frametitle}{%
  \nointerlineskip
  \begin{beamercolorbox}[wd=\textwidth, ht=2.1ex, dp=0.9ex,
                          leftskip=0.35cm, rightskip=0.35cm]{frametitle}%
    \usebeamerfont{frametitle}\insertframetitle%
  \end{beamercolorbox}%
  \vspace*{-0.4pt}%
  {\color{PyJaune}\rule{\textwidth}{1.6pt}}%
}

% ================= Pied de page =================
\setbeamerfont{footline}{size=\tiny}
\setbeamercolor{footline}{bg=white, fg=PyGris}
\newcommand{\PiedDePage}{Python en 24 heures \textbar\ Séance 02 : Boucles et algorithmes}
\setbeamertemplate{footline}{%
  \leavevmode%
  \hbox{%
    \begin{beamercolorbox}[wd=\paperwidth, ht=2.2ex, dp=1.1ex,
                            leftskip=0.35cm, rightskip=0.35cm]{footline}%
      \usebeamerfont{footline}\PiedDePage
      \hfill\insertframenumber\,/\,\inserttotalframenumber%
    \end{beamercolorbox}}%
}

\setbeamertemplate{navigation symbols}{}
\setbeamertemplate{itemize item}{\color{PyBleuMoyen}$\blacktriangleright$}
\setbeamertemplate{itemize subitem}{\color{PyBleuMoyen}\textbullet}

% ================= Blocs pedagogiques =================
\setbeamercolor{block title}{bg=PyBleu, fg=white}
\setbeamercolor{block body}{bg=PyBleuClair, fg=PyTexte}
\setbeamercolor{block title alerted}{bg=PyRouge, fg=white}
\setbeamercolor{block body alerted}{bg=PyRouge!9, fg=PyTexte}
\setbeamercolor{block title example}{bg=PyJaune, fg=PyNuit}
\setbeamercolor{block body example}{bg=PyJaune!14, fg=PyTexte}

% ================= Metadonnees =================
\hypersetup{
  pdftitle={Python en 24 heures : Séance 02 : Boucles et algorithmes},
  pdfauthor={HOUNKPE Mawulikplimi Roland et DJABON Ounimborbitibou},
  pdfsubject={De l'algorithmique à la programmation orientée objet},
  pdfkeywords={Python, algorithmique, POO, formation, exercices corrigés}
}

\begin{document}

{%
\setbeamercolor{background canvas}{bg=PyNuit}
\begin{frame}[plain]
\vspace{10mm}
{\color{PyJaune}\small\bfseries FORMATION « PYTHON EN 24 HEURES » \quad\textbar\quad SÉANCE 02\par}
\vspace{4mm}
{\color{white}\Huge\bfseries Répéter sans se répéter : boucles et algorithmes\par}
\vspace{4mm}
{\color{PyBleuClair}\large Chapitres 3 et 4 du Tome 1\par}
\vspace{7mm}
{\color{PyJaune}\rule{34mm}{2.2pt}\par}
\vspace{7mm}
{\color{white}\normalsize 1 h de cours (présentation Beamer) \quad$\bullet$\quad 1 h de pratique (notebook Jupyter)\par}
\vspace{2.5mm}
{\color{PyBleuClair}\normalsize HOUNKPE Mawulikplimi Roland et DJABON Ounimborbitibou\par}
\end{frame}
}
\begin{frame}{Objectifs de la séance}
\begin{itemize}
  \item Choisir entre la boucle TANT QUE et la boucle POUR selon la situation
  \item Utiliser un compteur et un accumulateur pour compter et additionner
  \item Écrire des boucles imbriquées pour traiter des tableaux à deux dimensions
  \item Implémenter la recherche du minimum, du maximum et le tri par sélection
  \item Identifier les deux pièges majeurs : la boucle infinie et le décalage de un
\end{itemize}
\end{frame}
\begin{frame}{Rappel de la séance 01}
Avant de répéter, revoyons les briques que vous maîtrisez déjà : elles seront utilisées à l'intérieur de chaque boucle.

\begin{itemize}
  \item Variables et affectation : une boîte étiquetée qui garde une valeur
  \item ENTRER et AFFICHER : dialoguer avec l'utilisateur
  \item Structures SI, SINON, SINON SI : choisir un chemin
  \item Opérateurs de test et connecteurs ET, OU, NON
  \item Méthode : analyser, décrire en pseudo-code, tester à la main
\end{itemize}
\end{frame}
\begin{frame}{Pourquoi répéter ?}
\begin{itemize}
  \item Afficher la table de 7 : dix instructions identiques ? Inacceptable
  \item Additionner 200 notes à la main : source d'erreurs garantie
  \item La boucle répète un bloc d'instructions \textbf{sans recopier le code}
  \item Deux outils selon la situation : \textbf{TANT QUE} (tant qu'il faut) et \textbf{POUR} (un nombre connu de fois)
\end{itemize}
\end{frame}

\begin{frame}[fragile]{La boucle TANT QUE}

\begin{lstlisting}[style=py]
ALGO CompteARebours
compteur reçoit 5
TANT QUE compteur > 0 FAIRE
    AFFICHER compteur
    compteur reçoit compteur moins 1
FIN TANT QUE
AFFICHER "Décollage !"
\end{lstlisting}
\vspace{0.5pt}

{\small\color{PyGris}\itshape Trois ingrédients : initialisation avant la boucle, condition de poursuite, mise à jour à l'intérieur.}

\end{frame}

\begin{frame}[fragile]{La boucle POUR}

\begin{lstlisting}[style=py]
ALGO TableDe7
POUR k ALLANT DE 1 A 10 FAIRE
    AFFICHER k, " fois 7 = ", k fois 7
FIN POUR
\end{lstlisting}
\vspace{0.5pt}

{\small\color{PyGris}\itshape Quand le nombre de répétitions est connu d'avance, POUR est plus clair que TANT QUE.}

\end{frame}

\begin{frame}{Choisir la bonne boucle}
\begin{columns}[t]
\begin{column}{0.48\textwidth}
{\bfseries\color{PyBleu}POUR}

\begin{itemize}
  \setlength{\itemsep}{1.5pt}\setlength{\parskip}{0pt}
  \item Nombre de tours \textbf{connu} à l'avance
  \item Parcours d'une suite ou d'un tableau
  \item Le compteur est géré automatiquement
\end{itemize}
\end{column}
\hspace{0.015\textwidth}
\begin{column}{0.48\textwidth}
{\bfseries\color{PyBleu}TANT QUE}

\begin{itemize}
  \setlength{\itemsep}{1.5pt}\setlength{\parskip}{0pt}
  \item Nombre de tours \textbf{inconnu} à l'avance
  \item Arrêt dicté par une condition
  \item Attention : la condition doit finir par devenir fausse
\end{itemize}
\end{column}
\end{columns}
\end{frame}

\begin{frame}[fragile]{Compteur et accumulateur}

\begin{lstlisting}[style=py]
ALGO MoyenneDe10Notes
somme reçoit 0
POUR i ALLANT DE 1 A 10 FAIRE
    ENTRER note
    somme reçoit somme + note
FIN POUR
moyenne reçoit somme / 10
AFFICHER "Moyenne : ", moyenne
\end{lstlisting}
\vspace{0.5pt}

{\small\color{PyGris}\itshape L'accumulateur `somme` démarre à 0 et grossit d'un terme à chaque tour : schéma à retenir toute l'année.}

\end{frame}

\begin{frame}[fragile]{Minimum et maximum d'une série}
\begin{columns}[t]
\begin{column}{0.42\textwidth}
\begin{itemize}
  \setlength{\itemsep}{1.5pt}\setlength{\parskip}{0pt}
  \item Idée : garder le champion courant dans une variable
  \item Au premier élément, le champion est ce premier élément
  \item Chaque nouvel élément remplace le champion s'il est plus petit (ou plus grand)
  \item Ne jamais initialiser à 0 : une série de notes négatives donnerait un faux maximum
\end{itemize}
\end{column}
\hspace{0.02\textwidth}
\begin{column}{0.55\textwidth}
\begin{lstlisting}[style=pysmall]
ALGO MaximumDeLaSerie
ENTRER n          # combien de valeurs
ENTRER premier
maxi reçoit premier
POUR i DE 2 A n FAIRE
    ENTRER valeur
    SI valeur > maxi ALORS
        maxi reçoit valeur
    FIN SI
FIN POUR
AFFICHER "Maximum : ", maxi
\end{lstlisting}
\end{column}
\end{columns}
\end{frame}

\begin{frame}[fragile]{Boucles imbriquées}

\begin{lstlisting}[style=py]
ALGO GrilleDeProduits
POUR ligne ALLANT DE 1 A 3 FAIRE
    POUR colonne ALLANT DE 1 A 4 FAIRE
        AFFICHER ligne, colonne
    FIN POUR
FIN POUR
\end{lstlisting}
\vspace{0.5pt}

{\small\color{PyGris}\itshape Pour chaque tour de la boucle externe, la boucle interne exécute sa propre boucle complète : 3 fois 4 tours au total.}

\end{frame}

\begin{frame}{Recherche dans un tableau}
\begin{itemize}
  \item Problème : la valeur v figure t elle dans la liste ?
  \item Parcourir les éléments un à un et \textbf{sortir dès la première rencontre}
  \item Si la boucle se termine sans trouver, la réponse est : absent
  \item Complexité : au pire n comparaisons pour n éléments
\end{itemize}
\end{frame}

\begin{frame}[fragile]{Recherche avec témoin}

\begin{lstlisting}[style=py]
ALGO RechercheDeV
trouve reçoit faux
POUR i ALLANT DE 1 A n FAIRE
    SI tableau[i] = v ALORS
        trouve reçoit vrai
    FIN SI
FIN POUR
SI trouve = vrai ALORS
    AFFICHER "Présent"
SINON
    AFFICHER "Absent"
FIN SI
\end{lstlisting}
\vspace{0.5pt}

{\small\color{PyGris}\itshape Le témoin `trouve` évite d'afficher plusieurs fois le même résultat et reste simple à traduire.}

\end{frame}

\begin{frame}[fragile]{Tri par sélection : l'idée}

\begin{lstlisting}[style=py]
ALGO TriParSelection
POUR i ALLANT DE 1 A n moins 1 FAIRE
    posMini reçoit i
    POUR j ALLANT DE i + 1 A n FAIRE
        SI tableau[j] < tableau[posMini] ALORS
            posMini reçoit j
        FIN SI
    FIN POUR
    échanger tableau[i] et tableau[posMini]
FIN POUR
\end{lstlisting}
\vspace{0.5pt}

{\small\color{PyGris}\itshape On cherche le plus petit élément restant et on le place à sa position définitive, tour après tour.}

\end{frame}

\begin{frame}{Attention}
\begin{alertblock}{Les deux pièges classiques}
La boucle infinie : si la variable de la condition n'est jamais mise à jour, le TANT QUE ne s'arrête jamais. Le décalage de un : de 1 à n fait n tours, de 0 à n en fait n + 1 ; décidez toujours explicitement des bornes. Dans les deux cas, exécutez la boucle à la main sur un petit exemple avec trois tours.
\end{alertblock}
\end{frame}

\begin{frame}{À retenir}
\begin{block}{L'essentiel de la séance 02}
POUR quand on compte, TANT QUE quand on attend un événement. Le schéma accumulateur (somme qui démarre à 0) et le schéma champion courant (maximum initialisé au premier élément) reviennent dans tous les algorithmes. Le tri par sélection combine boucles imbriquées et échange : apprenez à le dérouler à la main, c'est le meilleur moyen de comprendre toutes les boucles.
\end{block}
\end{frame}

\begin{frame}{Travaux pratiques (1 h) : Boucles et algorithmes classiques}
Cinq exercices pour muscler vos boucles : comptage, accumulation, test de primalité et un tri complet. Écrivez d'abord la logique en pseudo-code pour les exercices 3 à 5, puis codez.

\begin{itemize}
  \setlength{\itemsep}{1.5pt}\setlength{\parskip}{0pt}
  \item \textbf{Exercice 1 (Basique)} : Table de multiplication
  \item \textbf{Exercice 2 (Basique)} : Somme et moyenne d'une série
  \item \textbf{Exercice 3 (Intermédiaire)} : Factorielle
  \item \textbf{Exercice 4 (Intermédiaire)} : Nombre premier
  \item \textbf{Exercice 5 (Défi)} : Tri par sélection pas à pas
\end{itemize}
\vspace{4pt}

{\small\color{PyGris}Ouvrez le notebook « Seance\_02\_boucles\_et\_algorithmes.ipynb » pour commencer.}
\end{frame}
\begin{frame}{Pour la prochaine séance}
\begin{exampleblock}{Conseil de travail}
Séance 03 : premiers pas concrets en Python. Nous installerons nos habitudes d'interpréteur, découvrirons les types int, float, str et bool, les conversions et les f-strings. Refaites la table de 7 avec un print par ligne pour vous échauffer.
\end{exampleblock}
\end{frame}

\end{document}
